% Direct-authority diplomatic French for printed p.141.
\SourcePage{141}

\LayerTarget{tohoku:remark:satellites-existence:p2}
\noindent que pour tout foncteur covariant additif $F$ de $\mathbf{C}$ dans $\mathbf{C}'$, il
existe un $\partial$-foncteur universel défini en tous les degrés et se
réduisant à $F$ pour le degré zéro (i. e. tous les satellites $S^iF$
existent) il suffit manifestement de prouver que $S^1F$ et $S^{-1}F$ existent,
en vertu des formules
$S^1(S^iF)=S^{i+1}F$ si $i\geqq0$, $S^{-1}(S^iF)=S^{i-1}F$ si $i\leqq0$,
(dont la vérification à partir de la définition est triviale).
D'ailleurs, la recherche de $S^1F$ et $S^{-1}F$ sont deux problèmes duals, ces
deux foncteurs se permutent si on remplace $\mathbf{C}$ et $\mathbf{C}'$ par les catégories
duales.

\LayerTarget{tohoku:def:effaceable-coeffaceable}
Un foncteur additif $F$ de $\mathbf{C}$ dans $\mathbf{C}'$ est dit \textit{effaçable} si pour
tout $A\in \mathbf{C}$, on peut trouver un monomorphisme $u:A\to M$ tel que
$F(u)=0$ ; si $\mathbf{C}$ est tel que tout objet $A\in \mathbf{C}$ admet un effacement injectif
(cf. 1.10 remarque 1) il revient au même de dire que $F(u)=0$ pour tout
effacement injectif $u$ ; si enfin $\mathbf{C}$ est tel que tout objet $A\in \mathbf{C}$ admet
un monomorphisme dans un objet injectif $M$ (cf. 1.10.1) il revient au même de
dire que $F(M)=0$ pour tout objet injectif $M$. Dualement, $F$ est dit
\textit{coeffaçable} si pour tout $A\in \mathbf{C}$, on peut trouver un épimorphisme
$u:P\to A$ tel que $F(u)=0$.

\LayerTarget{tohoku:prop:2.2.1}
\indent\textsc{Proposition 2.2.1.} \textit{Soient $\mathbf{C},\mathbf{C}'$ deux catégories
abéliennes, $T=(T^i)\quad 0<i<b$ un $\partial$-foncteur exact (covariant ou
contravariant) de $\mathbf{C}$ dans $\mathbf{C}'$, avec $a>0$ Si $T^i$ est effaçable pour
$i>0$ alors $T$ est un $\partial$-foncteur universel, et la réciproque est
vraie si $\mathbf{C}$ est tel que tout objet $A\in \mathbf{C}$ admet un effacement injectif
\textup{(}\textit{cf.}\textup{ 1.10)}.}

\LayerTarget{tohoku:prop:2.2.1:proof:p1}
Il suffit pour la partie directe, de prouver par exemple que si
$(T'^0,T'^1)$ est un $\partial$-foncteur défini pour les degrés 0, 1 et $f^0$
un morphisme fonctoriel $T^0\to T'^0$, alors il existe un
morphisme $(f^0,f^1)$ unique de $(T^0,T^1)$ dans $(T'^0,T'^1)$ se réduisant à
$f^0$ pour le degré 0 (nous avons supposé $T$ covariant pour fixer les idées).
Soit $A\in \mathbf{C}$, considérons une suite exacte
$0\to A\to M\to A'\to 0$
telle que le morphisme $T^1(A)\to T^1(M)$ soit nul. Si on a pu
construire $f^1$, on aura un diagramme commutatif

\LayerTarget{tohoku:diagram:prop-2.2.1}
\[
\begin{array}{lclclcl}
T^0(M)&\to&T^0(A')&\to&T^1(A)&\to&T^1(M)\\
\big\downarrow&&\big\downarrow&&\big\downarrow&&\\
T'^0(M)&\to&T'^0(A')&\to&T'^1(A).&&
\end{array}
\]

\LayerTarget{tohoku:prop:2.2.1:proof:p2}
\noindent La première ligne étant exacte, on en conclut que le morphisme
$T^0(A')\to T^1(A)$ est surjectif, et par suite que le morphisme
$f^1(A):T^1(A)\to T'^1(A)$ est complètement déterminé par passage
au quotient à partie de
$f^0(A'):T^0(A)\to T'^0(A')$, ce qui prouve l'unicité de $f^1(A)$.
D'ailleurs, le diagramme précédent sans $f^1(A)$, compte tenu du fait que le
produit des deux morphismes de la deuxième ligne est nul, permet de définir un
morphisme $T^1(A)\to T'^1(A)$ de façon unique par la condition que
le diagramme reste commutatif. Des raisonnements standarts montrent que le
morphisme ainsi défini ne dépend pas du choix particulier de la suite exacte
$0\to A\to M\to A'\to 0$,
puis le fait que ce morphisme est fonctoriel, et
\textquotedblleft permute à $\partial$\textquotedblright. Cela prouve donc la
première partie de la proposition. La deuxième partie est contenue dans le
théorème d'existence suivant :

\LayerTarget{tohoku:thm:2.2.2:p1}
\indent\textsc{Théorème 2.2.2.} \textit{Soit $\mathbf{C}$ une catégorie
abélienne telle que tout objet $A\in \mathbf{C}$ admette un effacement injectif
\textup{(}\textit{cf.}\textup{ 1.10)}. Alors pour tout foncteur additif covariant $F$ sur $\mathbf{C}$, les
satellites $S^iF$ ($i\geqq0$) existent et sont des foncteurs effaçables pour
$i>0$. Pour que le $\partial$-foncteur universel $(S^iF)_{i\geqq0}$ soit exact,
il}

\TohokuContinuation{142}
